% =====================================================================
%  chunks/rev_app.tex —— 附录 读书笔记（tr-jia 撰）
% =====================================================================
\chapterreview{%
附录是 Novikov 为 1995 年英译本增写的一章，不是第 4 章的续编，而是「补记」：正文写于 1983--84 年、1986 年出俄文版，此后十年 3 维流形与纽结的纯拓扑涌现出一批新思想，故补此附录。因此它的语气与正文迥异——正文是教科书式的系统综述，附录更像写给同行的快报，先交代同期拓扑学的其他前沿（Floer 同调、$SO_2$-作用与椭圆亏格、拓扑量子场论、矩阵模型与瞬子），并且坦白作者对量子场论能否成为「流形上的严格分析」仍持保留态度。四节依次为：A.1 导言；A.2 纽结，从 Alexander 多项式的经典途径到 Jones 型多项式；A.3 Vassiliev 不变量；A.4 3 维流形的新不变量与拓扑量子场论。

主线是一个困境与三条出路。困境是：3 维流形的分类在「几何化」之后仍缺可计算的不变量，纽结的情形同样——虽已知道纽结由带标号的基本群决定，判定算法却至今未决。出路之一是多项式（A.2）：先用定向图上消去一个交叉的关系 $zP_{D_{n-1}}=P_{D_{n,+}}-P_{D_{n,-}}$ 与两条补充条件定义 Alexander 多项式，再以 $xH_{D_{n-1}}=yH_{D_{n,+}}-y^{-1}H_{D_{n,-}}$ 定义 HOMFLY，令 $x=y^{1/2}-y^{-1/2}$ 得 Jones 多项式 $J(y)$，统称 Jones 型多项式；Kauffman 给出非定向图的另一条途径（状态多项式 $Q$、括号多项式 $S$，$J(A^4,A^{-4})=(-A)^{-3w(D)}S(A,B;d)$）。出路之二是辫群（A.2 后半）：经 Markov 变换与 Markov 迹、Yang--Baxter 方程的解 $R$，把不变量写成 $A^{-w(a)}\Tr[\rho_n(a)]$；量子群是这一切的系统化，Hecke、Temperley--Lieb 与 Birman--Wenzl 代数则是其代数骨架；这条线的收尾是 Tait 的交错纽结——不可约交错图的交叉数是不变量（Kauffman--Murasugi）。出路之三是 Vassiliev 不变量（A.3）：把多项式不变量的可加性解释为对奇点求导，$f_{k-1}(G_{k-1}^+)-f_{k-1}(G_{k-1}^-)=f_k(G_k)$，第 $(k+1)$ 阶导数为零者称 $k$-不变量；弦图加两条关系得商群 $B_k$，$B=\sum_{k\ge0}B_k$ 在连通和下成 Hopf 环；Kontsevich 用多重积分造出 $Z(K)$，令 $\gamma=Z(U)$，则 $\gamma^{-s}Z(K)$ 是 Vassiliev 不变量。最后 A.4 收束到 TQFT：Witten 把 Reidemeister--Ray--Singer 挠的 Schwarz 表示推广为 Chern--Simons 型泛函积分 $Z(M^3)=\int DA\,e^{ikS(A)}$，严格化由 Reshetikhin--Turaev 完成——$M^3_L$ 由 $S^3$ 中带框链环 $L$ 作手术得到，Kirby 两个基本运算刻画其自由度，$TR(M^3_L)=\delta^{-i(L)}\langle\langle L\rangle\rangle$。

应记住的编号不多：附录只给两条编号公式，即 2 维 TQFT 的结合性方程 (4.1) $\eta^{ks}\partial_i\partial_j\partial_kF(t)=C^s_{ij}(t)$ 与 (4.2) $\eta_{ij}=\partial_1\partial_i\partial_jF(t)$；其余关键式皆无编号，只能照原书描述引用——Alexander 与 HOMFLY 的两条消去关系、Vassiliev 导数定义式、$TR$ 不变量与 Kontsevich 公式。Dubrovin 的 Frobenius 几何把这些收拢为一句话：联络 $\partial_i-zC_{ik}^j$ 的曲率为零，正是结合性方程成立的表现。

辨析有四处。其一，Kirby 运算是整节的技术枢纽：它把「3 维流形 = $S^3$ 上带框链环作手术」变成组合可算的对象，第二个运算里 $L_i\to L_i^1$、$L_j\to\bar L_j$ 的替换与「$2\pi$ 扭转」等价，宜对照图 A.13 读。其二，TQFT 的公理系统是 Segal--Atiyah 式的函子化表述（$V\mapsto H_V$、配边 $\mapsto$ 线性映射、$H_V=H_{\bar V}^*$），与物理学家的路径积分不是一回事，切勿把 $Z(M^3)=\int DA\,e^{ikS}$ 直接当公理。其三，原书 p.292 写出的 Chern--Simons「局部形式」$S(A)=\Tr(\partial_iA_i-\partial_jA_j+\frac23A^3)$ 不是坐标不变的写法（下标并未缩并），标准形式是 $\int_M\Tr(A\wedge dA+\frac23A\wedge A\wedge A)$，不要当标准式背。其四，「量子群」名不副实：作者明说它更像「离散化」而非「量子化」（p.287），并老实承认 Vassiliev 不变量除 Alexander 与 HOMFLY 的系数外尚无具体例子、与拓扑问题的联系还很浅（p.290）。书中笔误已在译注标出：Calaby--Yau（p.274、p.296，通行拼法 Calabi--Yau）、Hecke 关系 $\sigma_i-\sigma_{i-1}$ 应为 $\sigma_i-\sigma_i^{-1}$、$J(t,t^{-1})$、$V=S_1$ 应为 $S^1$、$T^2\subset M^2$ 等，正文照录，读时按批注校正。

评注：本附录最值得学的是组织方式：先给分类问题的困境，再给三条构造路线，最后清楚划线——已解决者（Reshetikhin--Turaev 的严格定义、Kontsevich 定理、Witten 猜想）与仍空着者（$TR$ 不变量尚未解决任何拓扑问题、Chern--Simons 的严格基础、Frobenius 几何的公理化）一目了然。这种「问题—路线—边界」的写法比文献综述有用。
}

% CHUNK rev_app DONE
